
\documentclass{article}
\usepackage{colortbl}
\usepackage{makecell}
\usepackage{multirow}
\usepackage{supertabular}

\begin{document}

\newcounter{utterance}

\centering \large Interaction Transcript for game `cladder', experiment `full\_v1.5\_default', episode 420 with score{-}Qwen3.5{-}9B{-}sp.
\vspace{24pt}

{ \footnotesize  \setcounter{utterance}{1}
\setlength{\tabcolsep}{0pt}
\begin{supertabular}{c@{$\;$}|p{.15\linewidth}@{}p{.15\linewidth}p{.15\linewidth}p{.15\linewidth}p{.15\linewidth}p{.15\linewidth}}
    \# & \multicolumn{2}{c}{Player} && \multicolumn{2}{c}{Game Master} \\
    \hline

    \theutterance \stepcounter{utterance}  
    & & & \multicolumn{4}{p{0.6\linewidth}}{
        \cellcolor[rgb]{0.9,0.9,0.9}{
            \makecell[{{p{\linewidth}}}]{
                \texttt{\tiny{[P1$\langle$GM]}}
                \texttt{You are an expert in causal inference. The following question is not a typical commonsense query, but rather a meticulously designed question created by a professor specializing in causal inference, intended to assess the students' mastery of the course content.} \\
\texttt{Imagine a self{-}contained, hypothetical world with only the following conditions, and without any unmentioned factors or causal relationships: Vaccination status has a direct effect on vaccination reaction and getting smallpox. Getting smallpox has a direct effect on smallpox survival. Vaccination reaction has a direct effect on smallpox survival. Method 1: We look directly at how vaccination status correlates with smallpox survival in general. Method 2: We look at this correlation case by case according to vaccination reaction. To understand how vaccination status affects smallpox survival, is it more correct to use the Method 1 than Method 2?} \\
\texttt{Start your answer with 'Yes' or 'No', followed by additional reasoning or evidence to support your explanation.} \\
            }
        }
    }
     \\ \\

    \theutterance \stepcounter{utterance}  
    & \multicolumn{4}{p{0.6\linewidth}}{
        \cellcolor[rgb]{0.561,0.737,0.561}{
            \makecell[{{p{\linewidth}}}]{
                \texttt{\tiny{[P1$\rangle$GM]}}
                \texttt{No} \\
\\ 
\texttt{**Reasoning:**} \\
\\ 
\texttt{In this hypothetical world, the causal structure can be represented as a Directed Acyclic Graph (DAG) where:} \\
\texttt{1.  **Vaccination Status** \$\textbackslash{}rightarrow\$ **Vaccination Reaction**} \\
\texttt{2.  **Vaccination Status** \$\textbackslash{}rightarrow\$ **Smallpox** (Getting smallpox)} \\
\texttt{3.  **Vaccination Reaction** \$\textbackslash{}rightarrow\$ **Smallpox Survival**} \\
\texttt{4.  **Smallpox** \$\textbackslash{}rightarrow\$ **Smallpox Survival**} \\
\\ 
\texttt{**Analysis of Method 1 (Total Effect):**} \\
\texttt{Method 1 examines the raw correlation between Vaccination Status and Smallpox Survival. This captures the **total causal effect**, which includes both the direct path (Vaccination Status \$\textbackslash{}rightarrow\$ Smallpox \$\textbackslash{}rightarrow\$ Survival) and the indirect path mediated through Vaccination Reaction (Vaccination Status \$\textbackslash{}rightarrow\$ Reaction \$\textbackslash{}rightarrow\$ Survival). In causal inference, when the question asks "how vaccination status affects smallpox survival" without specifying a conditional context, it refers to the total effect. Therefore, Method 1 correctly estimates the overall impact of the intervention (vaccination) on the outcome.} \\
\\ 
\texttt{**Analysis of Method 2 (Stratification/Conditioning):**} \\
\texttt{Method 2 conditions on **Vaccination Reaction**. In this DAG, **Vaccination Reaction** is a **mediator** (an intermediate variable) on the causal pathway from Vaccination} \\
            }
        }
    }
    & & \\ \\

    \theutterance \stepcounter{utterance}  
    & & & \multicolumn{2}{p{0.3\linewidth}}{
        \cellcolor[rgb]{0.9,0.9,0.9}{
            \makecell[{{p{\linewidth}}}]{
                \texttt{\tiny{[GM$|$GM]}}
                \texttt{No} \\
            }
        }
    }
    & & \\ \\

    \theutterance \stepcounter{utterance}  
    & & & \multicolumn{2}{p{0.3\linewidth}}{
        \cellcolor[rgb]{0.9,0.9,0.9}{
            \makecell[{{p{\linewidth}}}]{
                \texttt{\tiny{[GM$|$GM]}}
                \texttt{yes} \\
            }
        }
    }
    & & \\ \\

    \theutterance \stepcounter{utterance}  
    & & & \multicolumn{2}{p{0.3\linewidth}}{
        \cellcolor[rgb]{0.9,0.9,0.9}{
            \makecell[{{p{\linewidth}}}]{
                \texttt{\tiny{[GM$|$GM]}}
                \texttt{game\_result = LOSE} \\
            }
        }
    }
    & & \\ \\

\end{supertabular}
}

\end{document}
